\documentclass[../BTL.tex]{subfiles}
\begin{document}

% ═════════════════════════════════════════════════════════
%  PHẦN I — PHÂN TÍCH VÀ THIẾT KẾ HỆ THỐNG
% ═════════════════════════════════════════════════════════

\section{Quy trình tổng thể}
\label{section:3.1}

Hệ thống \textbf{CineRec} được xây dựng theo quy trình năm giai đoạn tuần tự,
đảm bảo tính logic và khả năng tái lập:

\begin{enumerate}[itemsep=0.4em]
    \item \textbf{Nghiên cứu lý thuyết (Giai đoạn 1):} Tìm hiểu tổng quan về
    \textit{Recommender Systems}, \textit{Collaborative Filtering},
    \textit{Matrix Factorization} và các chỉ số đánh giá (xem Chương~2).

    \item \textbf{Thu thập và tiền xử lý dữ liệu (Giai đoạn 2):} Lựa chọn
    dataset, làm sạch dữ liệu, phân tích khám phá (EDA) và xây dựng ma trận
    \textit{user--item}.

    \item \textbf{Xây dựng mô hình (Giai đoạn 3):} Triển khai ba mô hình CF,
    huấn luyện bằng \textit{Stochastic Gradient Descent}, chia tập
    \textit{train/test}.

    \item \textbf{Đánh giá và tối ưu (Giai đoạn 4):} Tính RMSE, MAE,
    Precision@K, Recall@K. Điều chỉnh siêu tham số.

    \item \textbf{Xây dựng demo (Giai đoạn 5):} Xây dựng API backend và giao
    diện web, kiểm tra toàn hệ thống.
\end{enumerate}

% ─────────────────────────────────────────────────────────
\section{Thu thập và tiền xử lý dữ liệu}
\label{section:3.2}

\subsection{Nguồn dữ liệu}
\label{section:3.2.1}

Dữ liệu được lấy từ bộ dữ liệu \texttt{MovieLens ml-100k}
(\url{https://grouplens.org/datasets/movielens/100k/})~\cite{harper2015movielens}.
Bộ dữ liệu gồm hai file gốc chính:

\begin{itemize}[itemsep=0.3em]
    \item \texttt{u.data}: Chứa $100.000$ dòng dữ liệu \textit{rating} theo
    định dạng \textit{tab-separated}: \texttt{userId}, \texttt{movieId},
    \texttt{rating}, \texttt{timestamp}.

    \item \texttt{u.item}: Chứa thông tin $1.682$ phim theo định dạng
    \textit{pipe-separated}: \texttt{movieId}, \texttt{title},
    \texttt{release\_date}, \texttt{imdb\_url} và $19$ cột thể loại nhị phân.
\end{itemize}

Hai file trên được parse thành \texttt{ratings.csv} và \texttt{movies.csv}
để thuận tiện cho việc đọc bằng Pandas.

\subsection{Tiền xử lý}
\label{section:3.2.2}

\begin{description}[leftmargin=1.5em, itemsep=0.5em]
    \item[\textbf{Đọc và chuẩn hóa:}] Đọc \texttt{u.data} với
    \textit{encoding} \texttt{latin-1}, tách $19$ cột thể loại nhị phân thành
    chuỗi \texttt{genres}.

    \item[\textbf{Kiểm tra dữ liệu:}] Xác nhận không có giá trị thiếu (NaN)
    trong cột \texttt{rating}. Kiểm tra phân phối \textit{rating} (nguyên,
    thuộc thang $1$–$5$).

    \item[\textbf{Xây dựng ma trận \textit{user-item}:}] Tạo \textit{pivot
    table} ($943 \times 1.615$) với \texttt{userId} làm \textit{index},
    \texttt{movieId} làm cột. Các ô chưa có \textit{rating} điền bằng $0$.
    Kích thước cột là $1.615$ (nhỏ hơn $1.682$) vì tập \textit{train} chỉ
    chứa phim xuất hiện trong $80\%$ \textit{rating} đầu tiên của mỗi người
    dùng.

    \item[\textbf{Chia tập \textit{train/test}:}] Áp dụng
    \textit{temporal split} (chi tiết tại mục~\ref{section:3.4}).
\end{description}

\subsection{Phân tích khám phá dữ liệu (EDA)}
\label{section:3.2.3}

\subsubsection*{Phân phối rating}

Hình~\ref{fig:rating_dist} thể hiện phân phối các mức \textit{rating}.
\textit{Rating} $4$ chiếm tỷ lệ cao nhất (${\approx}28\%$), tiếp theo là
\textit{rating} $3$ (${\approx}27\%$) và $5$ (${\approx}21\%$).
\textit{Rating} $1$ ít nhất (${\approx}6\%$). Phân phối lệch trái — người
dùng có xu hướng đánh giá cao, phản ánh \textit{positive bias} điển hình
trong dữ liệu \textit{rating} tự nguyện.

\begin{figure}[H]
\centering
\includegraphics[width=0.82\textwidth]{phan_phoi_rating}
\caption{Phân phối \textit{rating} trên bộ dữ liệu MovieLens ml-100k
($n = 100.000$). Phân phối lệch trái với \textit{rating} $4$ chiếm tỷ lệ
cao nhất.}
\label{fig:rating_dist}
\end{figure}

\subsubsection*{Hoạt động người dùng và độ phổ biến phim}

Phân phối số \textit{rating} theo người dùng rất lệch
(\textit{long-tail distribution}): một số ``super users'' đánh giá hơn $700$
phim, trong khi nhiều người dùng chỉ đánh giá $20$ phim (ngưỡng tối thiểu).
Tương tự, một số phim kinh điển được đánh giá hàng trăm lần
(\textit{Star Wars}: $584$ lượt). Hai đặc điểm \textit{long-tail} này là
nguyên nhân trực tiếp dẫn đến độ thưa thớt ${\approx}93{,}70\%$ của ma trận
\textit{user-item}.

% ─────────────────────────────────────────────────────────
\section{Thiết kế kiến trúc hệ thống}
\label{section:3.3}

\subsection{Kiến trúc tổng thể}
\label{section:3.3.1}

Hệ thống được thiết kế theo mô hình \textit{Client-Server} với ba lớp
(Hình~\ref{fig:architecture}):

\begin{description}[leftmargin=1.5em, itemsep=0.5em]
    \item[\textbf{Lớp dữ liệu} (\textit{Data Layer}):] File CSV được parse từ
    \texttt{u.data} và \texttt{u.item}. Ma trận \textit{user-item} được xây
    dựng trong bộ nhớ RAM khi khởi động.

    \item[\textbf{Lớp mô hình} (\textit{Model Layer}):] Ba class Python:
    \texttt{UserBasedCF}, \texttt{ItemBasedCF}, \texttt{MatrixFactorization}.
    Tất cả được huấn luyện một lần khi khởi động và lưu trong bộ nhớ để phục
    vụ dự đoán \textit{online} nhanh.

    \item[\textbf{Lớp trình bày} (\textit{Presentation Layer}):] FastAPI cung
    cấp RESTful API. Frontend HTML/CSS/JS gọi API và hiển thị kết quả.
\end{description}

\begin{figure}[H]
\centering
\scalebox{0.78}{%
\begin{tikzpicture}[
    node distance=0.5cm and 1.0cm,
    every node/.style={font=\small},
    layer/.style={draw=gray!50, fill=#1, rounded corners=6pt,
                  minimum width=3.2cm, minimum height=0.8cm,
                  align=center, text width=3.0cm},
    arrow/.style={-{Stealth[length=6pt]}, thick, #1},
    label/.style={font=\scriptsize\itshape, text=gray!70},
]
\node[layer=blue!10,   draw=blue!40]               (browser)
    {Web Browser\\{\scriptsize HTML/CSS/JS}};
\node[layer=orange!12, draw=orange!50,
      right=1.6cm of browser]                      (fastapi)
    {\texttt{FastAPI}\\{\scriptsize uvicorn :8000}};
\node[layer=green!12,  draw=green!50,
      above right=0.0cm and 1.5cm of fastapi]      (ubcf)
    {\texttt{UserBasedCF}\\{\scriptsize K=30}};
\node[layer=green!12,  draw=green!50,
      right=1.5cm of fastapi]                      (ibcf)
    {\texttt{ItemBasedCF}\\{\scriptsize K=30}};
\node[layer=green!12,  draw=green!50,
      below right=0.0cm and 1.5cm of fastapi]      (mf)
    {\texttt{MatrixFact.}\\{\scriptsize Funk SVD K=20}};
\node[layer=purple!10, draw=purple!40,
      right=1.5cm of ibcf]                         (matrix)
    {Ma trận $R$\\{\scriptsize $943{\times}1615$ (RAM)}};
\node[layer=purple!10, draw=purple!40,
      below=0.55cm of matrix]                      (csv)
    {\texttt{ratings.csv}\\{\scriptsize \texttt{movies.csv}}};

\draw[arrow=blue!60]   (browser.east) --
    node[above,label]{HTTP Request}  (fastapi.west);
\draw[arrow=blue!60]   (fastapi.west) --
    node[below,label]{JSON Response} (browser.east);
\draw[arrow=orange!70] (fastapi.east) |- (ubcf.west);
\draw[arrow=orange!70] (fastapi.east) -- (ibcf.west);
\draw[arrow=orange!70] (fastapi.east) |- (mf.west);
\draw[arrow=green!60]  (ubcf.east)   -| (matrix.north);
\draw[arrow=green!60]  (ibcf.east)   -- (matrix.west);
\draw[arrow=green!60]  (mf.east)     -| (matrix.south);
\draw[arrow=purple!60] (csv.north)   --
    node[right,label]{startup}       (matrix.south);

\begin{pgfonlayer}{background}
    \node[draw=orange!30, fill=orange!5, rounded corners=8pt,
          fit=(fastapi), inner sep=8pt,
          label={[font=\scriptsize\bfseries,text=orange!70]above:
                 Presentation Layer}] {};
    \node[draw=green!30,  fill=green!4,  rounded corners=8pt,
          fit=(ubcf)(ibcf)(mf), inner sep=8pt,
          label={[font=\scriptsize\bfseries,text=green!70]above:
                 Model Layer}] {};
    \node[draw=purple!30, fill=purple!4, rounded corners=8pt,
          fit=(matrix)(csv), inner sep=8pt,
          label={[font=\scriptsize\bfseries,text=purple!70]above:
                 Data Layer}] {};
\end{pgfonlayer}
\end{tikzpicture}
}% end scalebox
\caption{Kiến trúc ba lớp của hệ thống \textbf{CineRec} và luồng dữ liệu khi
xử lý một yêu cầu gợi ý}
\label{fig:architecture}
\end{figure}

\subsection{Thiết kế API}
\label{section:3.3.2}

\begin{table}[H]
    \centering
    \caption{Danh sách API endpoints}
    \label{tab:api}
    \begin{tabular}{llp{7cm}}
        \toprule
        \textbf{Endpoint} & \textbf{Method} & \textbf{Chức năng} \\
        \midrule
        \texttt{GET /}
            & GET & Trả về giao diện web (\texttt{index.html}) \\
        \texttt{GET /api/stats}
            & GET & Thống kê dataset + \textit{metrics} mô hình \\
        \texttt{GET /api/recommend/\{id\}}
            & GET & Gợi ý phim cho \texttt{user\_id}
                    (params: \texttt{method}, \texttt{n}) \\
        \texttt{GET /api/health}
            & GET & Kiểm tra trạng thái server \\
        \bottomrule
    \end{tabular}
\end{table}

\subsection{Cấu trúc dự án}
\label{section:3.3.3}

\begin{verbatim}
recsys_ml100k/
+-- data/
|   +-- u.data       +-- u.item
|   +-- ratings.csv  +-- movies.csv
+-- static/index.html
+-- models.py        (UserBasedCF, ItemBasedCF, MatrixFactorization)
+-- app.py           (FastAPI backend)
+-- run.py           (uvicorn launcher)
+-- demo_console.py  (terminal demo)
\end{verbatim}

% ─────────────────────────────────────────────────────────
\section{Chiến lược chia tập Train/Test}
\label{section:3.4}

Đồ án áp dụng \textbf{temporal split}: với mỗi người dùng, $80\%$
\textit{rating} sớm nhất làm \textit{train}, $20\%$ muộn nhất làm
\textit{test} (Hình~\ref{fig:temporal_split}). Chiến lược này phản ánh đúng
kịch bản thực tế và tránh \textit{data leakage}.

\begin{figure}[H]
\centering
\begin{tikzpicture}[every node/.style={font=\small}]
\draw[-{Stealth[length=7pt]}, thick, gray!70]
    (-0.3,0) -- (11.5,0)
    node[right, font=\small\itshape, text=gray!60] {Thời gian};
\node[font=\scriptsize, text=gray!60, below=4pt] at (0,  0) {Quá khứ};
\node[font=\scriptsize, text=gray!60, below=4pt] at (11, 0) {Hiện tại};

\foreach \y/\lbl in {0.55/1, -0.15/2, -0.95/{943}}{
    \fill[blue!35] (0,\y)     rectangle (8.8,\y+0.6);
    \fill[red!40]  (8.8,\y)   rectangle (11, \y+0.6);
    \draw[gray!50] (0,\y)     rectangle (11, \y+0.6);
    \node[font=\scriptsize\bfseries, text=white] at (4.4,  \y+0.3)
        {TRAIN (80\%)};
    \node[font=\scriptsize\bfseries, text=white] at (9.9,  \y+0.3)
        {TEST (20\%)};
    \node[left=4pt, font=\scriptsize] at (0,\y+0.3) {User $u_\lbl$};
}
\node[font=\small, text=gray!50] at (5.5,-0.55) {$\vdots$};
\node[left=4pt, font=\scriptsize, text=gray!50] at (0,-0.55) {$\vdots$};

\draw[dashed, thick, gray!60] (8.8,1.3) -- (8.8,-1.5)
    node[below, font=\scriptsize\itshape, text=gray!60]
    {Điểm phân cách (80\%)};

\fill[blue!35] (3.8,-1.85) rectangle (4.5,-1.6);
\node[right=2pt, font=\scriptsize] at (4.5,-1.725)
    {Tập \textit{train}};
\fill[red!40]  (7.2,-1.85) rectangle (7.9,-1.6);
\node[right=2pt, font=\scriptsize] at (7.9,-1.725)
    {Tập \textit{test}};
\end{tikzpicture}
\caption{\textit{Temporal split} theo từng người dùng: $80\%$
\textit{rating} sớm nhất làm \textit{train}, $20\%$ muộn nhất làm
\textit{test}. Mỗi người dùng được chia độc lập theo mốc thời gian của
chính họ.}
\label{fig:temporal_split}
\end{figure}

\begin{table}[H]
    \centering
    \caption{Kết quả chia tập train/test (temporal split 80/20)}
    \label{tab:split}
    \begin{tabular}{llll}
        \toprule
        \textbf{Tập} & \textbf{Số \textit{rating}} & \textbf{Tỷ lệ}
                     & \textbf{Mục đích} \\
        \midrule
        Train & ${\approx}80.000$ & ${\approx}80\%$
              & Huấn luyện mô hình \\
        Test  & ${\approx}20.000$ & ${\approx}20\%$
              & Đánh giá mô hình \\
        \bottomrule
    \end{tabular}
\end{table}

% ═════════════════════════════════════════════════════════

\end{document}
